\chapter{Einleitung}
\label{sec:einl}

Die Gliederung des Hauptteils dieses Dokuments ist nur als Platzhalter gedacht. 
Im Folgenden sind zwei Vorschl�ge f�r Gliederungen dargestellt, die
aber stets am konkreten Fall �berpr�ft und in der Regel angepasst werden
m�ssen.

Bitte beachten: Umfangreiche Programmlistings geh�ren nicht in die Arbeit, 
sondern auf eine beigef�gte CDROM. Programmbeispiele k�nnen auszugsweise
in der Arbeit gelistet werden, wenn sie zum Verst�ndnis notwendig sind.

\section{Literaturarbeit}
\label{sec:gliederung-lit}
\begin{enumerate}
  \item �berblick (oder: Zusammenfassung, "`Executive Summary"', alles Wichtige f�r den "`Manager"' oder Schnellleser)
  \item	Fragestellung (oder: Ziele, Ausgangspunkt, Motivation)
  \item	�bersicht �ber den Stand der Wissenschaft und Technik (Beschreibung der L�sungsans�tze, Beispiele etc. in einzelnen Abschnitten)
  \item	Bewertung der einzelnen untersuchten Ans�tze, Beispiele etc., Identifikation von Defiziten
  \item	Synthese: Erstellung einer Gesamtschau, allgemeine Prinzipien, Beschreibung einer eigenen Sicht auf das Problem, evtl. auch eigene Vorschl�ge
  \item	Zusammenfassung (Erkl�rung des Nutzens), Ausblick
\end{enumerate}
Anhang: eventuell recherchierte Texte, Produktbeschreibungen, etc.

\section{System- und Softwareentwicklungen}
\label{sec:gliederung-sys}
\begin{enumerate}
 \item	�berblick (oder: Zusammenfassung, "`Executive Summary"', alles Wichtige f�r den "`Manager"' oder Schnellleser)
 \item	Problemstellung (oder: Ziele, Ausgangspunkt), Vorgesehener Benutzerkreis, Bed�rfnisse der Benutzer
 \item	Stand der Technik (Wie wird das Problem bisher gel�st, wo sind die Defizite)
 \item	Gew�hlter L�sungsansatz (allgemeines Prinzip, welche Werkzeuge, z.B. Programmiersprachen werden verwendet)
 \item	Beschreibung der durchgef�hrten Arbeiten
 \item	Ergebnis (z.B. Screenshots mit Erl�uterungen, Bedienanleitung)
 \item	Zusammenfassung (Erkl�rung des Nutzens), Ausblick
\end{enumerate}
Anhang: evtl. (ausgew�hlte) Programmbeispiele.

